\ea\label{ex:text12-1}
O atovo eve o meho taba to pakupaku inu ni raara en Teapu.\\
\gll o	atovo	eve	o	meho	taba	to	paku{$\sim$}paku	inu	ni	ra=ara	en	teapu\\
     \textsc{art}3.\textsc{sg}	sago.palm.thatch	3\textsc{sg}.\textsc{pron}	\textsc{art}3.\textsc{sg}	other3.\textsc{sg}	thing	\textsc{rel}	\textsc{red}{$\sim$}make	house	\textsc{appl}	1\textsc{pl}.\textsc{in}.\textsc{ipfv}=1\textsc{pl}.\textsc{in}.\textsc{pron}	\textsc{dem}3	Teop\\
\glt `The sago palm thatch is something we use to build houses here in the Teop regions.'\footnote{lit. `The sago palm thatch, it (is) another thing that we make houses with here in Teop.'}
\z

\ea\label{ex:text12-2}
Eve to vavaakopo riori toon na inu tea varavihi kahi bona huan, a roava ae a iravata. A kiu vaa teo atovo he na paku vonaen:\\
\gll eve	to	va{$\sim$}vaa-kopo	ri=ori	toon	n=a	inu	tea	varavihi	kahi	bona	huan	a	roava	ae	a	iravata	a	kiu	vaa	te=o	atovo	he	na	paku	vo=na=en\\
     3\textsc{sg}.\textsc{pron}	\textsc{rel}	\textsc{red}{$\sim$}\textsc{caus}-lie.flat	3\textsc{pl}.\textsc{ipfv}=3\textsc{pl}.\textsc{pron}	top	3\textsc{sg}.\textsc{poss}=\textsc{art}2.\textsc{sg}	house	\textsc{compl}1	protect	from	\textsc{obj}.\textsc{art}2.\textsc{sg}	rain	\textsc{art}2.\textsc{sg}	sun	and1	\textsc{art}2.\textsc{sg}	storm	\textsc{art}2.\textsc{sg}	work	\textsc{lnk}	\textsc{prep}=\textsc{art}3.\textsc{sg}	sago.palm	\textsc{advers}	\textsc{tm}2	do	like=3\textsc{sg}.\textsc{ipfv}=\textsc{dem}3\\
\glt `It is put on top of the house to protect (it) against the rain, the sun and the storm.  The work on the sago palm is done as follows:'\footnote{lit. `It (is) what they put on top of the house ...'}
\z

\ea\label{ex:text12-3}
Beara suguna teo hum to tei nae a atovo, eara toro viana vamataa o vuaba vai o manoto vamataa, bara to vapepeana vamataa nana mee.\\
\gll be=ara	suguna	te=o	hum	to	tei	nae	a	atovo	e-ara	toro	viana	va-mataa	o	vuaba	vai	o	manoto	va-mataa	bara	to	va-pe{$\sim$}peana	va-mataa	nana	mee\\
     when=1\textsc{pl}.\textsc{in}.\textsc{pron}	arrive	\textsc{prep}=\textsc{art}3.\textsc{sg}	place	\textsc{rel}	stay	3\textsc{sg}.\textsc{ipfv}	\textsc{art}3.\textsc{pl}	sago.palm	\textsc{h}-1\textsc{pl}.\textsc{in}.\textsc{pron}	\textsc{tm}7	choose	\textsc{advr}-good	\textsc{art}3.\textsc{sg}	one	now	\textsc{art}3.\textsc{sg}	mature	\textsc{advr}-good	and2	\textsc{rel}	\textsc{mult}-\textsc{red}{$\sim$}have.branches	\textsc{advr}-good	3\textsc{sg}.\textsc{ipfv}	also\\
\glt `When we have arrived at the place where the sago palms are, we must carefully choose one that is fully mature and also has many branches.'
\z

\ea\label{ex:text12-4}
Beara viana vakavara, eara repaa tagava e ge eara repaa peha e.\\
\gll be=ara	viana	va-kavara	e-ara	re=paa	tagava	e	ge	e-ara	re=paa	peha	e\\
     when=1\textsc{pl}.\textsc{in}.\textsc{pron}	choose	\textsc{advr}-complete	\textsc{h}-1\textsc{pl}.\textsc{in}.\textsc{pron}	\textsc{consec}=\textsc{tm}3	fell	3\textsc{sg}.\textsc{pron}	or	\textsc{h}-1\textsc{pl}.\textsc{in}.\textsc{pron}	\textsc{consec}=\textsc{tm}3	climb	3\textsc{sg}.\textsc{pron}\\
\glt `When we have chosen (it), we fell it or we climb it.'
\z

\ea\label{ex:text12-5}
Tea peha, eara na peha rara, eara repaa mosi koa maa a peana gina beo atovo sun vakisi pahi nana.\\
\gll te=a	peha	e-ara	na	peha	rara	e-ara	re=paa	mosi	koa	maa	a	pea-na	gina	be=o	atovo	sun	vakisi	pahi	nana\\
     \textsc{prep}=\textsc{art}2.\textsc{sg}	climb	\textsc{h}-1\textsc{pl}.\textsc{in}.\textsc{pron}	\textsc{tm}2	climb	1\textsc{pl}.\textsc{in}.\textsc{ipfv}	\textsc{h}-1\textsc{pl}.\textsc{in}.\textsc{pron}	\textsc{consec}=\textsc{tm}3	cut	just	hither	\textsc{art}3.\textsc{pl}	branch-\textsc{derel}	on.the.top	when=\textsc{art}3.\textsc{sg}	sago.palm	stand	still	though	3\textsc{sg}.\textsc{ipfv}\\
\glt `As for the climbing, we climb it, and then we only cut the branches on the top, while the sago palm is still standing.'
\z

\ea\label{ex:text12-6}
Tea tagava, eara na tagava rara o atovo, ore obete nana putaa beara mosi rara a peana.\\
\gll te=a	tagava	e-ara	na	tagava	rara	o	atovo	o=re	obete	nana	putaa	be=ara	mosi	rara	a	pea-na\\
     \textsc{prep}=\textsc{art}2.\textsc{sg}	fell	\textsc{h}-1\textsc{pl}.\textsc{in}.\textsc{pron}	\textsc{tm}2	fell	1\textsc{pl}.\textsc{in}.\textsc{ipfv}	\textsc{art}3.\textsc{sg}	sago.palm	3.\textsc{pron}=\textsc{consec}	lie	3\textsc{sg}.\textsc{ipfv}	on.the.ground	when=1\textsc{pl}.\textsc{in}.\textsc{pron}	cut	1\textsc{pl}.\textsc{in}.\textsc{ipfv}	\textsc{art}3.\textsc{pl}	branch-\textsc{derel}\\
\glt `As for the felling, we fell the sago palm so that it lies on the ground, while we are cutting off the branches.'
\z

\ea\label{ex:text12-7}
Tea mosi peana he eara mee toro viana batara, are goe tea gono a peana vai to saka tara vamataa haana vai, a peana a tasitasiri ge to tara vakobusu nana.\\
\gll te=a	mosi	pea-na	he	e-ara	mee	toro	viana	bata=ra	a=re	goe	tea	gono	a	pea-na	vai	to	saka	tara	va-mataa	haa=na	vai	a	pea-na	a	tasi{$\sim$}ta-siri	ge	to	tara	va-kobusu	nana\\
     \textsc{prep}=\textsc{art}2.\textsc{sg}	cut	branch-\textsc{derel}	\textsc{advers}	\textsc{h}-1\textsc{pl}.\textsc{in}.\textsc{pron}	also	\textsc{tm}7	choose	\textsc{cont}=1\textsc{pl}.\textsc{in}.\textsc{ipfv}	1\textsc{pl}.\textsc{in}.\textsc{pron}=\textsc{consec}	must.not	\textsc{compl}1	get	\textsc{art}3.\textsc{pl}	branch-\textsc{derel}	\textsc{dem}6	\textsc{rel}	\textsc{neg}1	look	\textsc{advr}-good	\textsc{neg}2=3\textsc{sg}.\textsc{ipfv}	now	\textsc{art}3.\textsc{pl}	branch-\textsc{derel}	\textsc{art}3.\textsc{pl}	\textsc{red}{$\sim$}\textsc{acaus}-tear	or	\textsc{rel}	look	\textsc{advr}-moist	3\textsc{sg}.\textsc{ipfv}\\
\glt `As for the cutting of branches, we must also be selective, so that we won't get the branches that do not look good, branches that are torn or that look young and moist.'
\z

\ea\label{ex:text12-8}
Beara mosi vakavara a pea atovo, eara kou kosi vai a paka atovo kahi bona pea nae.\\
\gll be=ara	mosi	va-kavara	a	pea	atovo	e-ara	nao	vai	tea	kiu	vaa	tea	kosi	e-ara	kou	kosi	vai	a	paka	atovo	kahi	bona	pea	na=e\\
     when=1\textsc{pl}.\textsc{in}.\textsc{pron}	cut	\textsc{advr}-complete	\textsc{art}3.\textsc{pl}	branch	sago.palm	\textsc{h}-1\textsc{pl}.\textsc{in}.\textsc{pron}	go	now	\textsc{compl}1	work	\textsc{lnk}	\textsc{compl}1	cut.leaf	\textsc{h}-1\textsc{pl}.\textsc{in}.\textsc{pron}	\textsc{emph}2	cut.leaf	now	\textsc{art}3.\textsc{pl}	leaf	sago.palm	from	\textsc{obj}.\textsc{art}2.\textsc{sg}	branch	3\textsc{sg}.\textsc{poss}=3\textsc{sg}.\textsc{pron}\\
\glt `When we have cut the sago palm branches, we now go on to do the work of cutting off the leaves, we now cut the sago palm leaves off their branch.'
\z

\ea\label{ex:text12-9}
Be kavara o hum bona, eara nao vai tea kiu vaa tea varovaromo nia paka atovo.\\
\gll be	kavara	o	hum	bona	e-ara	nao	vai	tea	kiu	vaa	tea	varo{$\sim$}varomo	ni=a	paka	atovo\\
     when	finished	\textsc{art}3.\textsc{sg}	part	\textsc{ana}.\textsc{sg}	\textsc{h}-1\textsc{pl}.\textsc{in}.\textsc{pron}	go	now	\textsc{compl}1	work	\textsc{lnk}	\textsc{compl}1	\textsc{red}{$\sim$}stack.up	\textsc{appl}=\textsc{art}3.\textsc{pl}	leaf	sago.palm\\
\glt `When this part is finished, we now go on with the work of stacking up the sago palm leaves.'
\z

\ea\label{ex:text12-10}
Tea varovaromo he eara na varomo ni rara a ta taonim ge a ta peha-totoka paka atovo.\\
\gll te=a	varo{$\sim$}varomo	he	e-ara	na	varomo	ni	rara	a	ta	taonim	ge	a	ta	peha-totoka	paka	atovo\\
     \textsc{prep}=\textsc{art}2.\textsc{sg}	\textsc{red}{$\sim$}stack.up	\textsc{advers}	\textsc{h}-1\textsc{pl}.\textsc{in}.\textsc{pron}	\textsc{tm}2	stack.up	\textsc{appl}	1\textsc{pl}.\textsc{in}.\textsc{ipfv}	\textsc{art}2.\textsc{sg}	piece	five	or	\textsc{art}2.\textsc{sg}	piece	one-on.top	leaf	sago.palm\\
\glt `When stacking (them) up, we stack up five or six sago palm leaves.'
\z

\ea\label{ex:text12-11}
Bea kiu vaa tea varovaromo kavara, eara kou paku vai a peha dee i. Eara repaa kisi e, kisi rara bona kukan a kisi, eara repaa dee voosu.\\
\gll be=a	kiu	vaa	tea	varo{$\sim$}varomo	kavara	e-ara	kou	paku	vai	a	peha	dee	i	e-ara	re=paa	kisi	e	kisi	rara	bona	kukan	a	kisi	e-ara	re=paa	dee	voosu\\
     when=\textsc{art}2.\textsc{sg}	work	\textsc{lnk}	\textsc{compl}1	\textsc{red}{$\sim$}stack.up	finished	\textsc{h}-1\textsc{pl}.\textsc{in}.\textsc{pron}	\textsc{emph}2	make	now	\textsc{art}2.\textsc{sg}	one2.\textsc{sg}	load	\textsc{dem}1	\textsc{h}-1\textsc{pl}.\textsc{in}.\textsc{pron}	\textsc{consec}=\textsc{tm}3	tie	3\textsc{sg}.\textsc{pron}	tie	1\textsc{pl}.\textsc{in}.\textsc{ipfv}	\textsc{obj}.\textsc{art}2.\textsc{sg}	three	\textsc{art}2.\textsc{sg}	bundle	\textsc{h}-1\textsc{pl}.\textsc{in}.\textsc{pron}	\textsc{consec}=\textsc{tm}3	carry	go.home\\
\glt `When the work of stacking up is finished, we make one load. Then we tie it up by tying up three bundles, (and) then we carry (it) home.'
\z

\ea\label{ex:text12-12}
A peha aba na antee nana tea kosi bono dorana o peana\\
\gll a	peha	aba	na	antee	nana	tea	kosi	bono	dorana	o	pea-na\\
     \textsc{art}2.\textsc{sg}	one2.\textsc{sg}	person	\textsc{tm}2	can	3\textsc{sg}.\textsc{ipfv}	\textsc{compl}1	cut.leaves	\textsc{obj}.\textsc{art}3.\textsc{sg}	four	\textsc{art}3.\textsc{sg}	branch-\textsc{derel}\\
\glt `A single person can cut off the leaves of four branches.'
\z

\ea\label{ex:text12-13}
Ei beara suguna vaan, eara paku pete vai a meha kiu, a kiu vaa tea siri.\\
\gll ei	be=ara	suguna	vaan	e-ara	paku	pete	vai	a	meha	kiu	a	kiu	vaa	tea	siri\\
     \textsc{dem}2	when=1\textsc{pl}.\textsc{in}.\textsc{pron}	arrive	village	\textsc{h}-1\textsc{pl}.\textsc{in}.\textsc{pron}	do	subsequently	now	\textsc{art}2.\textsc{sg}	other2.\textsc{sg}	work	\textsc{art}2.\textsc{sg}	work	\textsc{lnk}	\textsc{compl}1	tear.off\\
\glt `Now, when we have arrived in the village, we do the next work, the work of tearing.'
\z

\ea\label{ex:text12-14}
Tea siri, eara na siri kahi batara a paka no atovo bona sikiri nae, ore tamee vamataa beara huha batara.\\
\gll te=a	siri	e-ara	na	siri	kahi	bata=ra	a	paka	n=o	atovo	bona	sikiri	na=e	o=re	tamee	va-mataa	be=ara	huha	bata=ra\\
     \textsc{prep}=\textsc{art}2.\textsc{sg}	tear.off	\textsc{h}-1\textsc{pl}.\textsc{in}.\textsc{pron}	\textsc{tm}2	tear.off	from	\textsc{cont}=1\textsc{pl}.\textsc{in}.\textsc{ipfv}	\textsc{art}2.\textsc{sg}	leaf	3\textsc{sg}.\textsc{poss}=\textsc{art}3.\textsc{sg}	sago.palm	\textsc{obj}.\textsc{art}2.\textsc{sg}	midrib	3\textsc{sg}.\textsc{poss}=3\textsc{sg}.\textsc{pron}	3.\textsc{pron}=\textsc{consec}	flexible	\textsc{advr}-good	when=1\textsc{pl}.\textsc{in}.\textsc{pron}	sew	\textsc{cont}=1\textsc{pl}.\textsc{in}.\textsc{ipfv}\\
\glt `As for the tearing, we tear the midrib off the sago palm leaf, so that it is flexible when we are sewing (the thatch sheets).'
\z

\ea\label{ex:text12-15}
Beara siri vakavara, eara repaa nao pete tea kiu vaa tea varovaromo vahabana nie, eara repaa kaakaku e.\\
\gll be=ara	siri	va-kavara	e-ara	re=paa	nao	pete	tea	kiu	vaa	tea	varo{$\sim$}varomo	vahabana	ni=e	e-ara	re=paa	kaa-kaku	e\\
     when=1\textsc{pl}.\textsc{in}.\textsc{pron}	tear.off	\textsc{advr}-complete	\textsc{h}-1\textsc{pl}.\textsc{in}.\textsc{pron}	\textsc{consec}=\textsc{tm}3	go	subsequently	\textsc{compl}1	work	\textsc{lnk}	\textsc{compl}1	\textsc{red}{$\sim$}stack.up	again	\textsc{appl}=3\textsc{sg}.\textsc{pron}	\textsc{h}-1\textsc{pl}.\textsc{in}.\textsc{pron}	\textsc{consec}=\textsc{tm}3	\textsc{red}{$\sim$}fold	3\textsc{sg}.\textsc{pron}\\
\glt `When we have finished the tearing, we go on with the work of stacking them (the leaves) up again, (and) then we fold them.'
\z

\ea\label{ex:text12-16}
Evehee eara toro kaku murina batara to peho ipinoo.\\
\gll evehee	e-ara	toro	kaku	murina	bata=ra	to	peho	ipinoo\\
     but	\textsc{h}-1\textsc{pl}.\textsc{in}.\textsc{pron}	\textsc{tm}7	fold	follow	\textsc{cont}=1\textsc{pl}.\textsc{in}.\textsc{ipfv}	\textsc{nspec}3.\textsc{sg}	one3.\textsc{sg}	measure\\
\glt `but we must fold (them) according to one measure.'\footnote{fig. `but we must fold (and) follow one measure.'}
\z

\ea\label{ex:text12-17}
Bea kaakaku kavara, eara repaa vaobete vatataunun e, eara repaa kisi vahabana nie.\\
\gll be=a	kaa-kaku	kavara	e-ara	re=paa	va-obete	va-ta{$\sim$}ta-unun	e	e-ara	re=paa	kisi	vahabana	ni=e\\
     when=\textsc{art}2.\textsc{sg}	\textsc{red}{$\sim$}fold	finished	\textsc{h}-1\textsc{pl}.\textsc{in}.\textsc{pron}	\textsc{consec}=\textsc{tm}3	\textsc{caus}-lie	\textsc{mult}-\textsc{red}{$\sim$}\textsc{acaus}-gather	3\textsc{sg}.\textsc{pron}	\textsc{h}-1\textsc{pl}.\textsc{in}.\textsc{pron}	\textsc{consec}=\textsc{tm}3	tie	again	\textsc{appl}=3\textsc{sg}.\textsc{pron}\\
\glt `When the folding is finished, we put them together, one on top of the other, then we tie them up again.'
\z

\newpage
\ea\label{ex:text12-18}
Eara repaa vaobete vamataa e teo hum to raka vamataa nana.\\
\gll e-ara	re=paa	va-obete	va-mataa	e	te=o	hum	to	raka	va-mataa	nana\\
     \textsc{h}-1\textsc{pl}.\textsc{in}.\textsc{pron}	\textsc{consec}=\textsc{tm}3	\textsc{caus}-lie	\textsc{advr}-good	3\textsc{sg}.\textsc{pron}	\textsc{prep}=\textsc{art}3.\textsc{sg}	place	\textsc{rel}	dry	\textsc{advr}-good	3\textsc{sg}.\textsc{ipfv}\\
\glt `We carefully lay them in a place that is completely dry.'
\z

\ea\label{ex:text12-19}
Pasi obete nom teebona, beara ihuana rara o bon vaa tea huha.\\
\gll pasi	obete	nom	tee-bona	be=ara	ihuana	rara	o	bon	vaa	tea	huha\\
     \textsc{tm}4	lie	\textsc{ipfv}	\textsc{prep}-\textsc{ana}.\textsc{sg}	when=1\textsc{pl}.\textsc{in}.\textsc{pron}	wait	1\textsc{pl}.\textsc{in}.\textsc{ipfv}	\textsc{art}3.\textsc{sg}	day	\textsc{lnk}	\textsc{compl}1	sew\\
\glt `They will lie there, while we are waiting for the sewing day.'
\z

\ea\label{ex:text12-20}
Eara na antee rara tea pahin huha bata ei bea tataraha matatopo nana ge ei bea naba kiu bero rori.\\
\gll e-ara	na	antee	rara	tea	pahin	huha	bata	ei	be=a	tataraha	matatopo	nana	ge	ei	be=a	naba	kiu	bero	rori\\
     \textsc{h}-1\textsc{pl}.\textsc{in}.\textsc{pron}	\textsc{tm}2	can	1\textsc{pl}.\textsc{in}.\textsc{ipfv}	\textsc{compl}1	directly	sew	\textsc{cont}	\textsc{dem}2	when2.\textsc{sg}=\textsc{art}	kururu.bamboo.strip	prepared	3\textsc{sg}.\textsc{ipfv}	or	\textsc{dem}2	when=\textsc{art}2.\textsc{sg}	team	work	many	3\textsc{pl}.\textsc{ipfv}\\
\glt `We can now directly sew (the thatch), when the kururu bamboo strip (on which the leaves are sewn) is prepared, or when the group of workers is big.'
\z

\ea\label{ex:text12-21}
Evehee a bero amaa vuri toro obete roho.\\
\gll evehee	a	bero	a=maa	vuri	toro	obete	roho\\
     but	\textsc{art}2.\textsc{sg}	many	\textsc{art}2.\textsc{sg}=\textsc{plm}	time	\textsc{tm}7	lie	before\\
\glt `But many days (the bundles) must lie (there) before (we can start).'
\z

\ea\label{ex:text12-22}
Tea kiu vaa tea huha he na paku vonaen:\\
\gll te=a	kiu	vaa	tea	huha	he	na	paku	vo=na=en\\
     \textsc{prep}=\textsc{art}2.\textsc{sg}	work	\textsc{lnk}	\textsc{compl}1	sew	\textsc{advers}	\textsc{tm}2	do	like=3\textsc{sg}.\textsc{ipfv}=\textsc{dem}3\\
\glt `As for the work of sewing, (it) is done as follows:'
\z

\newpage
\ea\label{ex:text12-23}
Eara na vaobete rara a tataraha toon no kasuana, eara repaa parata o paka atovo, ei kou na kaku pisi momohu ara.\\
\gll e-ara	na	va-obete	rara	a	tataraha	toon	n=o	kasuana	e-ara	re=paa	parata	o	paka	atovo	ei~kou	na	kaku	pisi	momohu	ara\\
     \textsc{h}-1\textsc{pl}.\textsc{in}.\textsc{pron}	\textsc{tm}2	\textsc{caus}-lie	1\textsc{pl}.\textsc{in}.\textsc{ipfv}	\textsc{art}2.\textsc{sg}	kururu.bamboo.strip	top	3\textsc{sg}.\textsc{poss}=\textsc{art}3.\textsc{sg}	ground	\textsc{h}-1\textsc{pl}.\textsc{in}.\textsc{pron}	\textsc{consec}=\textsc{tm}3	open	\textsc{art}3.\textsc{sg}	leaf	sago.palm	because	\textsc{tm}2	break	fold	first	1\textsc{pl}.\textsc{in}.\textsc{pron}\\
\glt `We put the kuru bamboo strip on the ground, then we open a sago.palm leaf, because we had folded it before.'
\z

\ea\label{ex:text12-24}
Erau eara repaa vaobete maa a tataraha komano paka atovo,\\
\gll erau	e-ara	re=paa	va-obete	maa	a	tataraha	koma-n=o	paka	atovo\\
     then	\textsc{h}-1\textsc{pl}.\textsc{in}.\textsc{pron}	\textsc{consec}=\textsc{tm}3	\textsc{caus}-lie	hither	\textsc{art}2.\textsc{sg}	kururu.bamboo.strip	inside-3\textsc{sg}.\textsc{poss}=\textsc{art}3.\textsc{sg}	leaf	sago.palm\\
\glt `Then we lay the kuru bamboo strip inside the (folded) sago palm leaf.'
\z

\ea\label{ex:text12-25}
Eara repaa vaakopo taono vahae bona.\\
\gll e-ara	re=paa	vaa-kopo	taono	vaha=e	bona\\
     \textsc{h}-1\textsc{pl}.\textsc{in}.\textsc{pron}	\textsc{consec}=\textsc{tm}3	\textsc{caus}-lie.flat	put.on.top	again=3\textsc{sg}.\textsc{pron}	4\textsc{sg}.\textsc{pron}\\
\glt `Then we cover it up again.'\footnote{lit. `Then we lay (it) flat by putting it on top of it.'}
\z

\ea\label{ex:text12-26}
Eara repaa huu nio koopu toon na ta atovo.\\
\gll e-ara	re=paa	huu	ni=o	koopu	toon	n=a	ta	atovo\\
     \textsc{h}-1\textsc{pl}.\textsc{in}.\textsc{pron}	\textsc{consec}=\textsc{tm}3	stick	\textsc{appl}=\textsc{art}3.\textsc{sg}	bamboo	top	3\textsc{sg}.\textsc{poss}=\textsc{art}2.\textsc{sg}	piece	sago.palm.thatch\\
\glt `Then we stick the koopu bamboo (needle) from above into the upper part of the (folded) sago palm thatch sheet.'
\z

\ea\label{ex:text12-27}
Eara repaa tahi bono sikiri vai to siri kahi ara bono atovo teo vuri siri atovo.\\
\gll e-ara	re=paa	tahi	bono	sikiri	vai	to	siri	kahi	ara	bono	atovo	te=o	vuri	siri	atovo\\
     \textsc{h}-1\textsc{pl}.\textsc{in}.\textsc{pron}	\textsc{consec}=\textsc{tm}3	insert	\textsc{obj}.\textsc{art}3.\textsc{sg}	midrib	\textsc{dem}6	\textsc{rel}	tear.off	from	1\textsc{pl}.\textsc{in}.\textsc{pron}	\textsc{obj}.\textsc{art}3.\textsc{sg}	sago.palm.leaf	\textsc{prep}=\textsc{art}3.\textsc{sg}	time	tear	sago.palm.leaf\\
\glt `Then we insert the midrib that we have torn off the sago palm leaf when doing the sago palm leaf tearing.'
\z

\ea\label{ex:text12-28}
Eara na tahi hopo rara, eara repaa tahi tavusu maa.\\
\gll e-ara	na	tahi	hopo	rara	e-ara	re=paa	tahi	tavusu	maa\\
     \textsc{h}-1\textsc{pl}.\textsc{in}.\textsc{pron}	\textsc{tm}2	push	enter	1\textsc{pl}.\textsc{in}.\textsc{ipfv}	\textsc{h}-1\textsc{pl}.\textsc{in}.\textsc{pron}	\textsc{consec}=\textsc{tm}3	push	come.out	hither\\
\glt `We push it into (the hole of the folded leaf), and then we push it out (of the folded leaf).'
\z

\ea\label{ex:text12-29}
eara repaa vaobete vahabana nao o meho paka, eara repaa huha bata naora.\\
\gll e-ara	re=paa	va-obete	vahabana	nao	o	meho	paka	e-ara	re=paa	huha	bata	nao=ra\\
     \textsc{h}-1\textsc{pl}.\textsc{in}.\textsc{pron}	\textsc{consec}=\textsc{tm}3	\textsc{caus}-lie	again	thither	\textsc{art}3.\textsc{sg}	other3.\textsc{sg}	leaf	\textsc{h}-1\textsc{pl}.\textsc{in}.\textsc{pron}	\textsc{consec}=\textsc{tm}3	sew	\textsc{cont}	thither=1\textsc{pl}.\textsc{in}.\textsc{ipfv}\\
\glt `Then we lay there another leaf and then go on sewing.'
\z

\ea\label{ex:text12-30}
Eara pasi paku bata koa vo naora teebona, ore tavusu komana to haha nae a ta tataraha.\\
\gll e-ara	pasi	paku	bata	koa	vo	nao=ra	tee-bona	o=re	tavusu	komana	to	haha	nae	a	ta	tataraha\\
     \textsc{h}-1\textsc{pl}.\textsc{in}.\textsc{pron}	\textsc{tm}4	make	\textsc{cont}	just	like	thither=1\textsc{pl}.\textsc{in}.\textsc{ipfv}	\textsc{prep}-\textsc{ana}.\textsc{sg}	3.\textsc{pron}=\textsc{consec}	come.out	self	\textsc{rel}	end	3\textsc{sg}.\textsc{ipfv}	\textsc{art}2.\textsc{sg}	piece	kururu.bamboo.strip\\
\glt `We will do it continuously just like this until the kururu bamboo strip is used up.'\footnote{lit. `We will do it continously just like this, then it comes out where the kururu bamboo strip ends.'}
\z

\ea\label{ex:text12-31}
Teo nin na vaapeha koana.\\
\gll te=o	nin	na	vaa-peha	koa-na\\
     \textsc{prep}=\textsc{art}3.\textsc{sg}	nin-reed	\textsc{tm}2	\textsc{caus}-one	just=3\textsc{sg}.\textsc{ipfv}\\
\glt `With a nin reed instead of a bamboo strip, it is just the same.'
\z

\ea\label{ex:text12-32}
Evehee ahiki kou to koopu, eara na tahi bata koa ni naora o nin,\\
\gll evehee	ahiki	kou	to	koopu	e-ara	na	tahi	bata	koa	ni	nao=ra	o	nin\\
     but	not.exist	\textsc{emph}2	\textsc{nspec}3.\textsc{sg}	bamboo	\textsc{h}-1\textsc{pl}.\textsc{in}.\textsc{pron}	\textsc{tm}2	insert	\textsc{cont}	just	\textsc{appl}	thither=1\textsc{pl}.\textsc{in}.\textsc{ipfv}	\textsc{art}3.\textsc{sg}	nin.reed\\
\glt `If there is not any koopu bamboo, we insert a nin reed.'
\z

\ea\label{ex:text12-33}
ei kou o tamee, bara eara na vavakakaniva rara o hum tahi.\\
\gll ei~kou	o	tamee	bara	e-ara	na	va{$\sim$}va-kakaniva	rara	o	hum	tahi\\
     because	\textsc{art}3.\textsc{sg}	flexible	and2	\textsc{h}-1\textsc{pl}.\textsc{in}.\textsc{pron}	\textsc{tm}2	\textsc{mult}-\textsc{caus}-sharp	1\textsc{pl}.\textsc{in}.\textsc{ipfv}	\textsc{art}3.\textsc{sg}	part	push\\
\glt `because it is flexible, and we carve a sharp pointed end that we push into the leaf.'\footnote{lit. `because it is flexible and we sharpen the pushing part'}
\z

\ea\label{ex:text12-34}
Be kavara a kiu vaa tea huha, eara repaa paku pete a kiu vaa tea vaave.\\
\gll be	kavara	a	kiu	vaa	tea	huha	e-ara	re=paa	paku	pete	a	kiu	vaa	tea	vaave\\
     when	finished	\textsc{art}2.\textsc{sg}	work	\textsc{lnk}	\textsc{compl}1	sew	\textsc{h}-1\textsc{pl}.\textsc{in}.\textsc{pron}	\textsc{consec}=\textsc{tm}3	do	subsequently	\textsc{art}2.\textsc{sg}	work	\textsc{lnk}	\textsc{compl}1	thatch\\
\glt `When the work of sewing is finished, we do the next work of putting up the thatch.'
\z

\ea\label{ex:text12-35}
Tea vaave, eara na paku roho rara a poka, o hum to sun bata raara beara vaave batara.\\
\gll te=a	vaave	e-ara	na	paku	roho	rara	a	poka	o	hum	to	sun	bata	ra=ara	be=ara	vaave	bata=ra\\
     \textsc{prep}=\textsc{art}2.\textsc{sg}	thatch	\textsc{h}-1\textsc{pl}.\textsc{in}.\textsc{pron}	\textsc{tm}2	make	before	1\textsc{pl}.\textsc{in}.\textsc{ipfv}	\textsc{art}2.\textsc{sg}	scaffold	\textsc{art}3.\textsc{sg}	place	\textsc{rel}	stand	\textsc{cont}	1\textsc{pl}.\textsc{in}.\textsc{ipfv}=1\textsc{pl}.\textsc{in}.\textsc{pron}	when=1\textsc{pl}.\textsc{in}.\textsc{pron}	thatch	\textsc{cont}=1\textsc{pl}.\textsc{in}.\textsc{ipfv}\\
\glt `Before putting up the thatch, we make a scaffold, a place on which we stand, when we are putting up the thatch.'
\z

\newpage
\ea\label{ex:text12-36}
Me tea vaave, eara na antee rara tea vaave hiava ge vaave gunaha.\\
\gll me	te=a	vaave	e-ara	na	antee	rara	tea	vaave	hiava	ge	vaave	gunaha\\
     and4	\textsc{prep}=\textsc{art}2.\textsc{sg}	thatch	\textsc{h}-1\textsc{pl}.\textsc{in}.\textsc{pron}	\textsc{tm}2	can	1\textsc{pl}.\textsc{in}.\textsc{ipfv}	\textsc{compl}1	thatch	go.up	or	thatch	go.down\\
\glt `And when thatching, we can thatch upwards (starting from the bottom) or thatch downwards (starting from the ridgepole).'
\z

\ea\label{ex:text12-37}
A buaku a aba pasi paku bata nom bona kiu vaa tea hee hiava ni bona atovo, bea meha papana vaave batari.\\
\gll a	buaku	a	aba	pasi	paku	bata	nom	bona	kiu	vaa	tea	hee	hiava	ni	bona	atovo	be=a	meha	papana	vaave	bata=ri\\
     \textsc{art}2.\textsc{sg}	two	\textsc{art}2.\textsc{sg}	person	\textsc{tm}4	do	\textsc{cont}	\textsc{ipfv}	\textsc{obj}.\textsc{art}2.\textsc{sg}	work	\textsc{lnk}	\textsc{compl}1	give	go.up	\textsc{appl}	\textsc{obj}.\textsc{art}3.\textsc{pl}	thatch.sheet	when=\textsc{art}2.\textsc{sg}	other2.sg	group	fix.thatch	\textsc{cont}=3\textsc{pl}.\textsc{ipfv}\\
\glt `Two people would do the work of reaching up the thatch sheets, while others would fix them.'
\z

\ea\label{ex:text12-38}
Eara na kisi ni rara o kara sareai, o nin ge o obai, teo varuvaru.\\
\gll e-ara	na	kisi	ni	rara	o	kara	sareai	o	nin	ge	o	obai	te=o	varuvaru\\
     \textsc{h}-1\textsc{pl}.\textsc{in}.\textsc{pron}	\textsc{tm}2	tie	\textsc{appl}	1\textsc{pl}.\textsc{in}.\textsc{ipfv}	\textsc{art}3.\textsc{sg}	string	sareai	\textsc{art}3.\textsc{sg}	nin	or	\textsc{art}3.\textsc{sg}	obai	\textsc{prep}=\textsc{art}3.\textsc{sg}	rafter\\
\glt `We tie (a thatch sheet) onto the rafter with the sareai vine, the nin reed or the calamus cane.'
\z

\ea\label{ex:text12-39}
Bea ta atovo obete maa tea varuvaru, eara repaa kisi taviri ie bona varuvaru.\\
\gll be=a	ta	atovo	obete	maa	te=a	varuvaru	e-ara	re=paa	kisi	taviri	e	bona	varuvaru\\
     when=\textsc{art}2.\textsc{sg}	sheet	sago.palm	lie	hither	\textsc{prep}=\textsc{art}3.\textsc{pl}	rafter	\textsc{h}-1\textsc{pl}.\textsc{in}.\textsc{pron}	\textsc{consec}=\textsc{tm}3	tie	go.around	3\textsc{sg}.\textsc{pron}	\textsc{obj}.\textsc{art}3.\textsc{pl}	rafter\\
\glt `When the thatch sheet lies on the rafters, we tie it (the string) around the rafters.'
\z

\newpage
\ea\label{ex:text12-40}
Be kavara, eara repaa dee hiava vahabana a ta meha ore suguna komana to antee nae a varuvaru.\\
\gll be	kavara	e-ara	re=paa	dee	hiava	vahabana	a	ta	meha	o=re	suguna	komana	to	antee	nae	a	varuvaru\\
     when	finished	\textsc{h}-1\textsc{pl}.\textsc{in}.\textsc{pron}	\textsc{consec}=\textsc{tm}3	reach	go.up	again	\textsc{art}2.\textsc{sg}	piece	other2.\textsc{sg}	3.\textsc{pron}=\textsc{consec}	arrive	indeed	\textsc{rel}	stop	3\textsc{sg}.\textsc{ipfv}	\textsc{art}3.\textsc{pl}	rafter\\
\glt `When it is finished, we reach up another (thatch sheet), until the end of the rafters is reached.'\footnote{lit. `... so that it arrives where the rafters stop.'}
\z

\ea\label{ex:text12-41}
Tea vaave hiava, eara na vaobete taono bata naora a ta atovo bona ta meha ore no suguna komana teo vatetekoio.\\
\gll te=a	vaave	hiava	e-ara	na	va-obete	taono	bata	nao=ra	a	ta	atovo	bona	ta	meha	o=re	no	suguna	komana	te=o	vatetekoio\\
     \textsc{prep}=\textsc{art}2.\textsc{sg}	thatch	go.up	\textsc{h}-1\textsc{pl}.\textsc{in}.\textsc{pron}	\textsc{tm}2	\textsc{caus}-lie	put.on.top.of	\textsc{cont}	thither=1\textsc{pl}.\textsc{in}.\textsc{ipfv}	\textsc{art}2.\textsc{sg}	piece	sago.palm	\textsc{obj}.\textsc{art}2.\textsc{sg}	piece	other2.\textsc{sg}	3.\textsc{pron}=\textsc{consec}	go	arrive	indeed	\textsc{prep}=\textsc{art}3.\textsc{sg}	ridgepole\\
\glt `When the thatching moves upwards, we put one sheet of sago palm thatch on top of the other, until the ridgepole is reached.'
\z

\ea\label{ex:text12-42}
Bea vaave gunaha, eara na taneo maara teo vatetekoio, ore gunaha maa beara hogo bata ni naora a ta atovo paana ta meha.\\
\gll be=a	vaave	gunaha	e-ara	na	taneo	maa=ra	te=o	vatetekoio	o=re	gunaha	maa	be=ara	hogo	bata	ni	nao=ra	a	ta	atovo	paa-n=a	ta	meha\\
     when=\textsc{art}2.\textsc{sg}	thatch	go.down	\textsc{h}-1\textsc{pl}.\textsc{in}.\textsc{pron}	\textsc{tm}2	start	hither=1\textsc{pl}.\textsc{in}.\textsc{ipfv}	\textsc{prep}=\textsc{art}3.\textsc{sg}	ridgepole	3.\textsc{pron}=\textsc{consec}	go.down	hither	when=1\textsc{pl}.\textsc{in}.\textsc{pron}	shove	\textsc{cont}	\textsc{appl}	thither=1\textsc{pl}.\textsc{in}.\textsc{ipfv}	\textsc{art}2.\textsc{sg}	piece	sago.palm	under-3\textsc{sg}.\textsc{poss}=\textsc{art}2.\textsc{sg}	piece	other2.\textsc{sg}\\
\glt `When the thatching moves downwards, we start at the ridgepole, and then it goes down,  when we shove one sago palm thatch sheet underneath the other.'
\z

\newpage
\ea\label{ex:text12-43}
Bea peha papa inu kavara, eara repaa taneo a meha papa inu.\\
\gll be=a	peha	papa	inu	kavara	e-ara	re=paa	taneo	a	meha	papa	inu\\
     when=\textsc{art}2.\textsc{sg}	one2.\textsc{sg}	side	house	finished	\textsc{h}-1\textsc{pl}.\textsc{in}.\textsc{pron}	\textsc{consec}=\textsc{tm}3	taneo	\textsc{art}2.\textsc{sg}	other2.\textsc{sg}	side	house\\
\glt `When one side of the house is finished, we start with the other side.'
\z

\ea\label{ex:text12-44}
Be kavara, ei kavara vai, eara repaa vaobete maa o kovuaka tea mahan vai to paa vatatana nom a vaave ei teo vatetekoio.\\
\gll be	kavara	ei	kavara	vai	e-ara	re=paa	va-obete	maa	o	kovuaka	te=a	mahan	vai	to	paa	vatatana	nom	a	vaave	ei	te=o	vatetekoio\\
     when	finished	\textsc{dem}2	finished	now	\textsc{h}-1\textsc{pl}.\textsc{in}.\textsc{pron}	\textsc{consec}=\textsc{tm}3	\textsc{caus}-lie	hither	\textsc{art}3.\textsc{sg}	top.thatch.sheet	\textsc{prep}=\textsc{art}2.\textsc{sg}	open.space	\textsc{dem}6	\textsc{rel}	\textsc{tm}3	meet	\textsc{ipfv}	\textsc{art}2.\textsc{sg}	thatching	\textsc{dem}2	\textsc{prep}=\textsc{art}3.\textsc{sg}	upper.ridge.pole\\
\glt `When (this other side) is finished, this (thatching) is finished now, and we put the top thatch sheet over this open space where the (two sides of the) thatch meet, this is (on top of) the upper ridge pole.'
\z
